\section{Concluding Remarks}
\label{sec:conclusion}

We surveyed the many active fronts of research in
multi-agent emergent language, empowered by advances in deep learning. Current simulations have become more realistic, running in 
setups which often include real-world images or grounded 3D environments,
giving rise to protocols with several intriguing properties. We conclude here by briefly listing a number of open questions and exciting directions for future work.

From a more theoretical perspective, we hope that more researchers from AI, linguistics and cognitive science will chip in 
with stronger hypotheses to shape experimental and analytical work, as
well as to provide insights into how to make computational models more
human-like.  Concepts playing an important role in the study of human communication and language, such as joint attention, theory
of mind and syntactic recursion are currently under-studied in the field of
multi-agent communication.  Hopefully, interdisciplinary insights will help modelers inject the right biases into agent
architectures and experimental setups. This will in turn result in simulations
that are not only more realistic in terms of investigating the roots of natural language, but also more
relevant to the kind of situations agents would encounter in actual human interaction.  Fortunately, toolkits are becoming available that facilitate
 entry into the area by scientists from other disciplines \cite<e.g.,>{Kharitonov:etal:2019}. 

While the studies presented here have already introduced a number of methods for
inspecting emergent protocols, important open issues remain with respect to how
to analyze emergent languages and whether it is possible to develop automated tools
that could speed up and generalize their analysis.

An even more serious issue that future research must address is that,
in the vast majority of current simulations, the
communication partners of agents during training and testing phases are the
same, i.e., we evaluate communication between agents that were trained
together. As such, it is not clear to what extent we are evaluating agents on
their ability to develop general communication skills, or simply their ability
to co-adapt to their learning environment and partners.

Perhaps the next big frontier with respect to
Artificial Intelligence lies in bringing the emergent language of agents to a
level of complexity and generality that will make them useful in applications.
However, just like human language did not probably emerge all at once,
we should outline precise desiderata for a minimally useful agent
\emph{proto-language}. At the same time, the huge progress currently being made
in corpus-based statistical natural language learning should be harnessed
to encourage the emergence of more interpretable and fluent protocols, thus
making multi-agent communication an integral part of human-centric AI. 


%We have seen that deep network communities can develop language-like
%communication protocols with several intriguing properties. Tools such
%as EGG (\cite{Kharitonov:etal:2019}) are now providing an easy
%entry-point into this budding field. It's time to think about how
%simulations can scale up, and become more relevant to both theoretical
%and applied work in cognitive science and AI. Stronger hypotheses from
%cognitive science could play a more direct role in shaping
%experimental and analytical work, as well as providing insights into
%how to make the computational models more human-like. The huge
%progress currently being made on corpus-based statistical natural
%language learning in AI could be harnessed to encourage the emergence
%of more interpretable and fluent communication protocols. Important
%issues are open both with respect to how to analyze emergent
%languages, and how to bring them to a level of complexity and
%generality that will make them useful in applications. We propose some
%exciting directions for future work in the Outstanding Questions box.
% In this survey, We tried to give a glimpse of the many active fronts of research in multi-deep-agent emergent language. We believe that some of the aspects\ldots


%\textbf{Alphago thing, GPT-2 ~\cite{Lewis:etal:2017} Hehe}

% \textbf{In all work presented in this manuscript,  the experimental setup dictates that the training and testing interaction partners of agents are stable, i.e., agents learn to communicate with a particular set of agents and all we test about is their co-adaptation skills, i.e., how well they overfitted to their partners. Thus it is arguable whether advances in the emergent communication correspond to advances in learning general communication skills transferable to other situations, including  different agent partners but also human partners. Recent experiments by Caroll et al. ~\cite{Carroll:etal:2019} in the context of humans co-operating with machines trained via self-play indicate that the transfer from agent-agent to agent-human is indeed a not a trivial one, even in the absence of language. Agents when trained with self-play tend to co-adapt to each other forming conventions that very often are idiosyncratic, since agents do not possess the same cognitive biases as humans (move 57 of alphaGo). While these experiments were reported without communication, we expect similar results would emerge also when co-operating using communication, as communication protocols tend to be uninterpretable.  Thus, moving forwards, it is important to guarantee that agents adopt protocols that allows them to communicate with humans. In fact, communicating via natural language should make transfer skills of general co-operation easier as it's a good universal representation.}

% \textbf{While language and cognitive science has already played a big role in emergent communication, given the fact that we are not there yet, we maybe need to search for more biases from cognitive process and human behaviour. Joint Attention (Tomasselo) is important. Humans model their interlocutors, theory of mind. While some work in emergent communication already goes in that direction, we need more. Especially once we will go towards this direction of cooperating with different partners than the training ones.}

% \textbf{NB:} the following ref is in the list we submitted with the
% proposal, but it was not referred to in our original draft:

% \begin{itemize}
% \item Das A, Kottur S, Moura JM, Lee S, Batra D. Learning cooperative visual dialog agents with deep reinforcement learning. In Proceedings of the IEEE International Conference on Computer Vision 2017.
% \end{itemize}

